\section*{Chapter \ref{ch:rc:stress-testing-and-scenarios} --- Stress Testing and Scenarios}

\begin{solution}{exo:rc:stress-testing-and-scenarios:1}
VaR is a quantile of a distribution estimated from a window; a stress scenario is a single severe
path from a different period, beyond the quantile and outside the window's regime. A correct 99\%
VaR is exceeded, by any amount, on 1\% of days.
\end{solution}

\begin{solution}{exo:rc:stress-testing-and-scenarios:2}
$\E[x_2\mid x_1] = \rho\frac{\sigma_2}{\sigma_1}x_1 = 0.6\times\frac{5}{10}\times40 = 12$ basis points.
\end{solution}

\begin{solution}{exo:rc:stress-testing-and-scenarios:3}
$10/2.5 = 4$ standard deviations.
\end{solution}

\begin{solution}{exo:rc:stress-testing-and-scenarios:4}
The two-year yield, correlated with the ten-year, rises too (80 basis points), and the book is short
two-year bonds: that leg gains and offsets part of the ten-year loss. A shock to one factor alone
ignores the hedges that the other factors' co-movement brings.
\end{solution}

\begin{solution}{exo:rc:stress-testing-and-scenarios:5}
The linear model ignores the short straddle's gamma: its most plausible scenario spreads the loss
over rates and currencies at 3.93 standard deviations, and in full revaluation loses USD~14.00
million, more than the target. The straddle loses on large currency moves in either direction, so a
nearer scenario, 2.47 standard deviations away, concentrated on EURUSD, reaches USD~10 million.
\end{solution}

\begin{solution}{exo:rc:stress-testing-and-scenarios:6}
That the moves would be the same in today's markets (levels, volatilities, correlations, the
structure of the market), that today's positions are held through the period without rebalancing,
and that the book could be liquidated at those prices.
\end{solution}

\begin{solution}{exo:rc:stress-testing-and-scenarios:7}
Without the straddle the linear distance is 3.84 standard deviations and the full-revaluation one
3.95: the bonds' positive convexity makes a loss slightly harder to reach than the linear model says,
the opposite of the straddle's effect.
\end{solution}

\begin{solution}{exo:rc:stress-testing-and-scenarios:8}
The probability uses a normal distribution and the recent covariance; real moves are fat-tailed and
correlations jump in crises, so four-sigma ten-day moves happen far more often (the historical
episodes of this chapter cost two to three times the ten-day VaR). The distance ranks scenarios; it
is not a probability.
\end{solution}

\begin{solution}{pb:rc:stress-testing-and-scenarios:1}
\textbf{1.} USD~1.5 billion.
\textbf{2.} USD~905 million.
\textbf{3.} While the client pays, the broker is hedged (it holds the stocks against the swaps); on
default the swaps vanish and the broker is left long the stocks, with the margin as its only cushion.
\textbf{4.} Almost nothing: the swaps and the hedges offset, and margin covers daily moves.
\textbf{5.} The client's total position across brokers: the same names, many times the size, whose
forced sale would move prices against every broker.
\textbf{6.} Falls of 8.7, 7.9, 7.1, 6.3 and 5.4\% in the five stocks, from the largest position to the
smallest.
\textbf{7.} 1.66 standard deviations; a normal probability of 4.9\% over five days.
\textbf{8.} The scenario loads on the direction $\Sigma\delta$: stocks with larger weights contribute more to
the loss per unit of move, and the common factor moves all of them together.
\textbf{9.} 7.5\% on every stock.
\textbf{10.} The distance rises to 4.42 standard deviations, a normal probability of $5.0\times10^{-6}$.
\textbf{11.} The positions were many days of trading volume; the broker competes with other brokers
selling the same names, and prices gap before any selling.
\textbf{12.} They add a common, correlated fall in exactly these stocks, which the covariance of
normal days does not contain: the scenario is both larger and likelier.
\textbf{13.} It was set low and static, while the concentration and volatility of the positions grew:
the cushion shrank relative to the risk.
\textbf{14.} The client's whole position (or at least its concentration in each name at this broker),
liquidity in days of volume, and volatility, re-set daily.
\textbf{15.} A limit on each name as a share of its daily volume and of the broker's capital, with
margin rising steeply with concentration.
\textbf{16.} The loss comes from a joint event (default and a concentrated fall) in a crowded
position, outside any window of normal days.
\textbf{17.} As a trigger for margin, limits and exits: when the distance falls under a threshold, act.
\textbf{18.} Its total exposure across brokers, its leverage, and its liquidity.
\textbf{19.} Named result: \emph{the family office from the broker's side}: falls of 8.7, 7.9, 7.1, 6.3 and
5.4\% over five days exhaust the 7.5\% margin, only 1.66 standard deviations away.
\textbf{20.} Which plausible worlds would break the firm, and how near they are.
\end{solution}

\begin{solution}{iq:rc:stress-testing-and-scenarios:1}
VaR: a probabilistic statement about ordinary losses, estimated from recent data, good for
day-to-day limits. \omterm{def:rc:stress-testing-and-scenarios:stress}{Stress test}: the loss in a specified severe scenario, no probability, good for
tail, concentration and regime changes. Use both.
\iqlookfor{complementary roles.}
\end{solution}

\begin{solution}{iq:rc:stress-testing-and-scenarios:2}
Specify the story's key shocks, fill the other factors by conditional expectation under a
(possibly stressed) covariance or a \omterm{def:rc:structural-credit-models:structural}{structural model}, check signs and magnitudes against history,
and review with the desks.
\iqlookfor{conditioning and expert review.}
\end{solution}

\begin{solution}{iq:rc:stress-testing-and-scenarios:3}
Minimise $\tfrac12x^\top\Sigma^{-1}x$ subject to $\delta^\top x = -L$: $x^* = -L\Sigma\delta/(\delta^\top\Sigma\delta)$, distance $L/\sqrt{\delta^\top\Sigma\delta}$.
\iqlookfor{the Lagrangian and the result.}
\end{solution}

\begin{solution}{iq:rc:stress-testing-and-scenarios:4}
An instantaneous, severe shock to many market factors applied to trading and fair-valued
positions of large banks, sized by the supervisor; VaR is a bank's own statistical estimate over a
day or ten.
\iqlookfor{supervisor-set, instantaneous, broad.}
\end{solution}

\begin{solution}{iq:rc:stress-testing-and-scenarios:5}
By client: default plus a concentrated fall of its largest positions over a realistic liquidation
horizon, with crowding across brokers; reverse stress to the margin; limits on concentration and
days of volume; dynamic margin.
\iqlookfor{default-conditional, concentration and liquidity.}
\end{solution}

\begin{solution}{iq:rc:stress-testing-and-scenarios:6}
Historical episodes by asset class and regime, \omterm{def:rc:stress-testing-and-scenarios:hypothetical}{hypothetical scenarios} for current risks, reverse
stress results, supervisory scenarios; versioned definitions, owners, review cycles, factor
mappings for new products, and reconciliation of results with VaR and P\&L.
\iqlookfor{content and governance.}
\end{solution}
